% ── Abstract ─────────────────────────────────────────────────────────────────
\begin{abstract}
Existing graph-based retrieval-augmented generation (RAG) systems apply a
single fixed retrieval strategy to every query regardless of its logical
structure.
We present \textbf{GRASP-RAG} (Graph Retrieval with Adaptive Structured
Paths), which introduces \emph{query-adaptive motif selection}: a lightweight
LLM classifier identifies the logical shape of each query---single-entity,
bridge, multi-entity, or broad---and routes it to the matching graph
traversal algorithm (ego-network, bidirectional shortest path, Steiner-tree
approximation, or LightRAG hybrid).
Each motif then structures its retrieved subgraph into a topology-aware
context that mirrors the answer's reasoning shape, so the LLM receives
evidence in the form the question requires.
When the classifier is uncertain, a confidence score triggers a graceful
fallback to LightRAG hybrid rather than a hard failure.
Evaluated on four domains of UltraDomain (Agriculture, CS, Legal, Health)
under reference-anchored multi-judge evaluation (GPT-4o-mini, Claude Haiku,
Gemini Flash; bias-corrected majority vote across five dimensions), GRASP-RAG
achieves average win rates of \textbf{73.4\% against LightRAG} and
\textbf{79.7\% against PathRAG}---the latter being consistent across all four
domains.
Classifier accuracy ranges from 40--72\% by domain, with bridge queries
classified reliably (87.7\% pooled) and broad queries misclassified most
often (10.8\%), yet per-type win rates remain positive for every query type,
confirming that the system degrades gracefully when classification is
imprecise.
A controlled ablation holding routing and retrieval fixed shows that
topology-aware context serialisation alone accounts for a 62.4\% win rate
over flat triple serialisation on Agriculture, confirming that structured
context presentation contributes independently of graph traversal.
\end{abstract}

\begin{IEEEkeywords}
retrieval-augmented generation, knowledge graph, graph RAG, multi-hop
reasoning, query classification, adaptive retrieval
\end{IEEEkeywords}
